\documentclass{article}

\usepackage{../../Style}

\title{Geometric Group Theory}
\author{Shing Tak Lam}
\date{Lent 2023\footnote{Based on lectures by Henry Wilton. Last updated \today.}}

\makeindex

\begin{document}

\maketitle

\tableofcontents

\lec

\section{Combinatorial group theory}

Combinatorial group theory is a sibling field to Geometric group theory. Both study infinite discrete groups.

\subsection{Free groups and presentations}

Let \(A = \{a_1, a_2, \dots\}\) be an \emph{alphabet}. A group \(F\) is \emph{free on \(A\)} if
\begin{enumerate}
    \item There is a map of sets \(A \to F\),
    \item for any group \(G\), and a map of sets \(A \to G\), there exists a unique group homomorphism \(F \to G\) such that
    % https://q.uiver.app/#q=WzAsMyxbMCwxLCJBIl0sWzEsMCwiRiJdLFsxLDEsIkciXSxbMCwxXSxbMSwyXSxbMCwyXV0=
    \[\begin{tikzcd}
        & F \\
        A & G
        \arrow[from=2-1, to=1-2]
        \arrow[from=1-2, to=2-2]
        \arrow[from=2-1, to=2-2]
    \end{tikzcd}\]
    commutes.
\end{enumerate}
This is a \emph{universal property}. As usual, this means that \(F\) is unique up to unique isomorphism. This shows that \(F\) is determined by \(A\), so we may write \(F = F(A)\).

However (as usual with definitions by universal property), we don't know if \(F(A)\) exists. We'll show this two different ways.

\textbf{1. Topologically:} Let
\[X = \bigvee_{a \in A}S^1\]
By the Seifert-van Kempen theorem,
\[\pi_1(X) \cong F(A)\]

\textbf{2. Combinatorially:} Let
\[A^* = \{\text{words in }A \sqcup A^{-1}\}\]
where \(A^{-1} = \{a_1^{-1}, \dots, a_n^{-1}\}\). For example,
\[1 = \emptyset, aa, aa^{-1}, aba^{-1}b^{-1}, a^{100}ba^{-100}b, \dots\]
A word is \emph{reducible} if it contains \(aa^{-1}\) or \(a^{-1}a\) as a subword for any \(a \in A\). Otherwise, it is \emph{reduced}. We can now define
\[F(A) = \left\{w \in A^* \text{ reduced}\right\}\]
The group operation is concatenation, followed by reduction. For example,
\[(abab^{-1})(b^2a) = abab^{-1}b^2a = ababa\]
Note that reduction terminates as each reduction step reduces the length. We won't check that this is well defined or associative. The identity element is \(1\), inverses is clear.

A \emph{presentation} consists of an alphabet \(A\), which we will call \emph{generators}, and a set \(R \sse F(A)\), which we will call \emph{relations}, and we write
\[\inner{A \mid R} = \inner{a_1, a_2, \dots \mid r_1, r_2, \dots}\]
This \emph{presents}  a group
\[G = \frac{F(A)}{\iinner{R}}\]
where \(\iinner{R}\) is the \emph{normal closure} of \(R\), i.e. the smallest normal subrgoup of \(F(A)\) containing \(R\).

\begin{example}
    Some examples of presentations which we have seen:
    \begin{align*}
        \inner{a \mid a^n} &\cong \Z/n\Z \\
        \inner{r, s \mid r^n, s^2, srsr} &\cong D_{2n} \\
        \inner{A\mid} &\cong F(A) \\
        \inner{a, b \mid aba^{-1}b^{-1}} &\cong \Z^2 \\
        \inner{a_1, \dots, a_g, b_1, \dots, b_g \mid a_1b_1a_1^{-1}b_1^{-1}\cdots a_gb_ga_g^{-1}b_g^{-1}} &\cong \pi_1(\Sigma_g)
    \end{align*}
    where \(\Sigma_g\) is the compact orientable surface of genus \(g\).
\end{example}
As we see, presentations arise when we write down fundamental groups of spaces. In fact, all groups arise this way.
\begin{corollary}
    [of Seifert-van Kampen]
    For
    \[G = \inner{a_1, a_2, \dots \mid r_1, r_2, \dots}\]
    there exists a space \(X\) with \(\pi_1(X) = G\).
\end{corollary}

\begin{proof}
    First, start of with a wedge of circles, one for each \(a_i\). Also consider a disjoint union of discs, one for each \(r_i\). Attach the \(i\)-cell along its boundary, which is a loop in the wedge of circles.
\end{proof}

This is called the \emph{presentation complex} of \(G\). For example, if we have \(\Z^2 = \inner{a, b \mid aba^{-1}b^{-1}}\), then \(X = \T^2\) is the \(2\)-torus.

In 1911, Max Dehn posed the following problems.
\begin{enumerate}
    \item (Word problem) Given \(w \in A^*\), determine whether or not \(w = 1\) in \(G = \inner{A \mid R}\). That is, whether or not \(w \in \iinner{R} \sse F(A)\).
    \item (Conjugacy problem) Given \(G = \inner{A \mid R}\), \(u, v \in A^*\), determine whether or not \(u\) is conjugate to \(v\) in \(G\).
    \item (Isomorphism problem) Given \(G = \inner{A \mid R}, H = \inner{B \mid S}\), determine whether or not \(G\) is isomorphic to \(H\).
\end{enumerate}

\lec

\begin{remark}
    The conjugacy problem is stronger than the word problem, since \(w = 1\) if and only if \(w\) is conjugate to \(1\).

    Dehn was motivated by topology, but the problems asks for algorithms. We will often solve them using geometry.

    All three were \emph{un}solved in the 1950s, as all three problems are algorithmically undecidable. Nevertheless, positive solutions are known for many reasonable classes of groups.
\end{remark}

Let \(A = \{a_1, \dots, a_n\}\) be a finite alphabet.

\begin{example}
    [word problem in free groups]

    Let \(w \in A^*\). If \(w\) is reduced, then \(w = 1\) if and only if \(w = \emptyset\). Otherwise, \(w\) contains a subword \(aa^{-1}\) for some \(a \in A \cup A^{-1}\). Cancelling \(aa^{-1}\) produces a word \(w' \in A^*\), such that \(w = w'\) in \(F(A)\), and \(\ell(w') = \ell(w) - 2 < \ell(w)\). This terminates after finitely many steps. 
\end{example}

We can also solve the conjugacy problem for free groups.

\begin{definition}
    There is a natural action of \(\Z\) on \(A^*\) permuting words. That is,
    \[1\cdot a_1\cdots a_k = a_2a_3\cdots a_ka_1\]
\end{definition}

The elements of \(\Z \cdot w\) are called the \emph{cyclic conjugates} of \(w\). Note that all \(u \in \Z w\) are conjugate to \(w\). The orbits \(\Z \backslash A^*\) are called \emph{cyclic words}. A word is \emph{cyclically reduced} if every cyclic conjugate is reduced.

\begin{example}
    \(aba^{-1}\) is reduced, but not cyclically reduced as \(ba^{-1}a\) is not reduced.
\end{example}

Note that if \(w\) is reduced, but not cyclically reduced, then
\[w = aw'a^{-1}\]
for some \(a \in A \cup A^{-1}\). Note that \(w'\) is both conjugate to \(w\) and shorter than \(w\). Hence after finitely many iterations, we can assume that \(w\) is cyclically reduced.

\begin{lemma}
    [conjugacy problem in free groups] If \(u, v \in F(A)\) are cyclically reduced, then \(u\) is conjugate to \(v\) if and only if the corresponding cyclic words are the same.
\end{lemma}

\begin{proof}
    If \(u, v\) have the same cyclic words, then \(v\) is a cyclic conjugate of \(u\), which we have seen is a conjugate of \(u\).

    Conversely, suppose \(u = gvg^{-1}\). By induction on \(\ell(g)\), we can assume \(g = a \in A \cup A^{-1}\). From this it follows that either \(v = a^{-1}v'\), or \(v = v'a\), as \(v\) is cyclically reduced. That is, \(u = v'a^{-1}\) or \(u = av'\). In both cases, they are cyclic conjugates.
\end{proof}

\subsection{Historical case study}

Let's briefly think about the state of topology in the early 20th century. Poincar\'e knew that homology classifies the compact two-dimensional surfaces. This motivated the

\begin{conjecture}
    [Poincar\'e conjecture, version 1] Let \(M\) be a compact connected 3-manifold, with
    \[\HH_*(M) = \begin{cases}
        \Z & * = 0, 3 \\
        0 & \text{otherwise}
    \end{cases}\]
    Then \(M\) is homeomorphic to \(S^3\).
\end{conjecture}

Such a 3-manifold \(M\) is called a \emph{homology sphere}.

\begin{theorem}
    [Poincar\'e, 1904] There exists a three dimensional homology sphere \(P\), with
    \[\pi_1(P) \surj A_5\]
\end{theorem}

The moral is that: homology is not enough, we need use \(\pi_1\) as well.

\begin{conjecture}
    [Poincar\'e conjecture, version 2] Let \(M\) be a compact connected \(3\)-manifold, with \(\pi_1(M) = 1\). Then \(M\) is homeomorphic to \(S^3\).
\end{conjecture}

This was proven by Perelman in 2003. Returning to the original conjecture, in 1910 Dehn wanted to construct more homology spheres.

\begin{theorem}
    [Dehn, 1910] There are infinitely many non-homeomorphic 3-dimensional homology spheres.
\end{theorem}

\begin{remark}
    The isomorphism problem is exactly what is needed to distinguish these manifolds.
\end{remark}

\lec

Here is Dehn's construction. Consider the trefoil knot \(K \sse \R^3 \sse S^3\). Let \(N(K)\) be a regular open neighbourhood of \(K\), i.e. a thickening of \(K\). Let \(N = S^3 \setminus N(K)\). This is a compact 3-manifold with boundart \(\partial N \cong \T^2\).

Computing,
\[\HH_*(N) = \begin{cases}
    \Z & * = 0, 1 \\
    0 & \text{otherwise}
\end{cases}\]
and
\[\pi_1(N) \cong \inner{x, y, z \mid x^2 = y^3 = z}\]
This follows from the fact that \(K\) is a torus knot, and we can compute the \(\pi_1\) of the complement of a torus knot. Note that the abelianisation map
\[\pi_1(N) \to \HH_1(N) \cong \Z\]
sends \(z\) to \(6\), \(x\) to \(3\) and \(y\) to \(2\). Moreover, the boundary torus \(T\) has \(\pi_1(T) = \Z^2\), generated by \(xy\) and \(z\). Under the abelianisation map, \(xy\) is sent to \(5\).

We can glue a solid torus \(D^2 \times S^1\) to \(N\), by a homeomorphism on the boundary. Let \(\lambda, \mu\) be the corresponding loops on \(\T^2\). The resulting manifold \(M_\phi = N \cup_\phi T\) is closed. \(\pi_1(M_\phi)\) depends on \(g = \phi_*(\mu)\). By Seifert-van Kapmen, 
\[\pi_1(M_\phi) = \frac{\pi_1(N)}{\iinner{g}}\]
Similarly,
\[\HH_1(M_\phi) = \frac{\Z}{\inner{[g] \in \HH^1(N)}}\]
To produce a homology sphere, we need to choose \(\phi\) such that \(g = \phi_*(\mu) = 1\) in \(\HH^1(N)\).

If \(g = (xy)^az^b\), then in \(\HH^1(N)\) this is mapped to \(5a + 6b \in \Z\). Choose \(a = 6n + 5\), and \(b = -(5n + 4)\) for \(n \in \Z\). He constructs \(\phi_n\) such that
\[\phi_n(\mu) = g_n = (xy)^{6n+5}z^{-(5n+4)}\]
His family of manifolds
\[D_n = N \cup_{\phi_n} U\]
has
\[\pi_1(D_n) = \inner{x, y, z \mid x^2 = y^2 = z, (xy)^{6n+5} = z^{5n + 4}}\]
The remaining challenge is to prove the groups \(G_n = \pi_1(D_n)\) for \(n \ge 0\) are pairwise non-isomorphic.

This is the isomorphism problem! In particular,
\[g_n = g_m \implies g_n \text{ and }g_m\text{ are conjugate} \implies G_n \cong G_m\]
So we also need to solve the word and conjugacy problem in \(\pi_1(N)\).

\subsection{van Kampen diagrams}

\begin{definition}
    A map \(f : Y \to X\) of cell complexes is called \emph{combinatorial} if for all \(k \in \Z_{\ge 0}\), and every \(k\)-cell \(e^k\) of \(Y\), \(f\) maps the interior \(\Int(e^k)\) homeomorphically to the interior of a \(k\)-cell of \(X\).
\end{definition}

Consider a presentation \(G = \inner{a_i \mid r_j}\), and the associated presentation complex \(X\).

\begin{definition}
    A \emph{(singular) disc diagram} is a compact contractible \(2\)-complex \(D\), with an embedding \(D \inj \R^2\). A disc diagram \(D\) is \emph{over \(X\)} if it is equipped with a combinatorial map \(D \to X\).
\end{definition}

\begin{example}
    Example of disc diagram.
\end{example}

Recall that \(X\) is given by a wedge of circles, with discs glued on for each relation. So the \(1\)-cells correspond to generators, \(2\)-cells go to relations (or cyclic conjugates, or inverses). With this:
\begin{itemize}
    \item every oriented \(1\)-cell of \(D\) is labelled with some \(a_i \in A\),
    \item so that each \(2\)-cell has boundary which is a cyclic conjugate of some \(r_j^{\pm 1}\).
\end{itemize}
Associated to each disc diagram \(D\), we have a \emph{boundary cycle}, which reads a (cyclic\footnote{or a word once we choose a base point.}) word \(w \in A^*\), which reduces to an element \(w' \in \iinner{r_j} \le F(A)\). To see this, \(D\) is contractible.

\(D\) is a van Kampen diagram for \(w\).



\end{document}